\begin{abstract}
\thispagestyle{kapitelkopfzeile}

% Optionale Kurzfassung Ihrer Arbeit (ca. eine halbe Seite).
% Inhalt: Problemstellung, Zielsetzung, Vorgehen und wichtigste Ergebnisse.
% Falls Sie keinen Abstract benötigen, können Sie den Inhalt dieser Datei
% leeren oder die Einbindung in der Hauptdatei auskommentieren.

\begin{otherlanguage}{ngerman}
\section*{Kurzfassung}

Die vorliegende Arbeit untersucht eine SAP-Fiori-Anwendung zur mobilen Fotodokumentation von Instandhaltungsmeldungen 
bei Woodward L'Orange, die bislang funktional, nicht jedoch systematisch hinsichtlich ihrer Gebrauchstauglichkeit geprüft wurde.
Die Forschungsfrage lautet, inwieweit die Anwendung die Interaktionsprinzipien der ISO~9241-110:2020 erfüllt und welche 
Optimierungspotenziale sich daraus ableiten lassen.
Aus der normeigenen Anwendungstabelle wurde ein Kriterienkatalog mit 41 im Nutzungskontext anwendbaren Prüffragen operationalisiert, 
gegen den die Anwendung in einer heuristischen Evaluation als Experteninspektion geprüft wurde.
Festgestellte Abweichungen wurden nach der Schweregradskala von Nielsen eingestuft.
Von den sieben Prinzipien sind vier erfüllt, zwei teilweise erfüllt und eines nicht erfüllt, 
wobei elf Abweichungen dokumentiert wurden.
Die Defizite betreffen vorrangig fehlende Rückmeldungen bei ungültigen Datei-Uploads sowie inhaltsleere 
und nicht handlungsleitende Fehlermeldungen.
Auf dieser Grundlage werden nach Schweregrad priorisierte Optimierungspotenziale abgeleitet.
\end{otherlanguage}

\begin{otherlanguage}{english}
\section*{Abstract}

This thesis examines an SAP Fiori application for the mobile photo documentation of maintenance requests 
at Woodward L'Orange, which has been tested for functionality but not systematically evaluated for usability.
The research question is to what extent the application complies with the interaction principles of ISO 9241-110:2020 and what 
opportunities for optimization can be derived from this.
A catalog of criteria comprising 41 test questions applicable to the usage context was operationalized from the standard’s application table, 
against which the application was evaluated in a heuristic evaluation conducted as an expert inspection.
Identified deviations were classified according to Nielsen’s severity scale.
Of the seven principles, four are met, two are partially met, and one is not met, 
with eleven deviations documented.
The shortcomings primarily concern a lack of feedback for invalid file uploads, as well as error messages that are devoid of content 
and do not guide user action.
Based on this, optimization opportunities are identified and prioritized according to severity.
\end{otherlanguage}
\end{abstract}
